\documentclass[a4paper]{article}
\usepackage{amsfonts}
\usepackage{amsmath}
\usepackage{amssymb}
\usepackage{graphicx}     

\begin{document}
  
\noindent \large Auteur: X Tenka Tenzin Choezin \\
\noindent \large Studierichting: Informatica\\
\noindent \large Oefeningenreeks 1E nr. 34\\

\medskip

\normalsize

Zoek een derde graad functie $A(z)$ met $\{1, -1\}$ als nulpunten.\\

$A(0) = 2 + i$ en $A(i) = 4$\\

$A(z) = az^3 + bz^2 + cz + R$\\

$A(0) = R = 2 + i$\\

$A(z) = az^3 + bz^2 + cz + 2 + i$\\

Vul 1 en -1 in $\Rightarrow \left\{
\begin{array}{l|l}
	a + b + c + 2 + i = 0  & 1\\
   	-c + b - c + 2 + i = 0 & 1
\end{array}
\right.$\\

$2b + 4 + 2i = 0$\\

$b = -2 - i$\\

Vul -1 en i in $\Rightarrow \left\{
\begin{array}{l|l}
	-a + b - c + 2 + i = 0 & 1\\
	-ai - b + ci + 2 + i = 4 & 1
\end{array}
\right.$\\

$-a - c - ai + ci = -2i$\\

$\left\{ 
\begin{array}{l}
	-a - c = 0 \\
	-ai + ci = -2i
\end{array}
\right.$\\

$\left\{ 
\begin{array}{l}
	a = -c \\
	ci + ci = -2i
\end{array}
\right.$\\

$\left\{ 
\begin{array}{l}
	a = -c \\
	2ci = -2i
\end{array}
\right. \Rightarrow a = 1 $ en $c = -1$\\

\textbf{Eindresultaat}\\

$A(z) = z^3 + (-2 - i)z^2 - z + 2 + i$\\

\begin{thebibliography}{999}
%\bibitem{a} Referentie
\end{thebibliography}
\end{document} 
